\pdfmapfile{=cm.map}
\pdfmapfile{+cmextra.map}
\pdfmapfile{+latxfont.map}
\pdfmapfile{+symbols.map}
\pdfmapfile{+euler.map}
\documentclass[11pt]{article}
\usepackage[margin=1in]{geometry}
\usepackage{amsmath,amssymb,amsthm,booktabs,hyperref,microtype}
\hypersetup{colorlinks=true,urlcolor=blue,linkcolor=blue}
\newtheorem{theorem}{Theorem}
\newtheorem{lemma}[theorem]{Lemma}
\newtheorem{proposition}[theorem]{Proposition}
\newcommand{\Q}{\mathbb Q}
\newcommand{\Z}{\mathbb Z}
\newcommand{\N}{\operatorname N}
\title{A smaller quartic for the H10 construction}
\author{Research and exact computations produced with ChatGPT\\for Gabriel Conant}
\date{November 6, 2026}
\begin{document}
\maketitle
\noindent This note records improvements to the constructions in
\href{https://github.com/elliotglazer/diophantine-z-q}{\texttt{elliotglazer/diophantine-z-q}}
and its open contribution branch. All variable counts include the input. The new result is a reduction of the defining polynomial, with the same rational zero projection as the specified source. It inherits that source's arithmetic dependencies; it is not an independent audit of the complete proof that the projection is $\Z$.

\section{Result and comparison}
The main repository was inspected at commit
\texttt{5856286a57fba19574a40704571ee70e79738281}.
Its displayed record is a quartic in 962 variables, with a fixed coefficient parameter $3^{2^{27}}$. The head of pull request~4 was fetched at
\texttt{772c4adf4a05e21255e4c1784e08e917fda4c68a}.
In addition to the packaged 933-variable polynomial, that branch contains a 909-variable direct-target candidate. The latter is the baseline for the strongest improvement here.

\begin{center}
\begin{tabular}{lrr}
\toprule
Construction & Total variables & Degree\\
\midrule
Main repository & 962 & 4\\
Contribution branch, ordinary coefficients & 909 & 4\\
Same branch, fixed symbolic coefficient & 892 & 4\\
New construction, ordinary coefficients & 850 & 4\\
New construction, fixed symbolic coefficient & 833 & 4\\
\bottomrule
\end{tabular}
\end{center}
Here ``ordinary coefficients'' means every coefficient is written as a decimal integer in the supplied table. The symbolic version uses the fixed integer
\[
 S=3^{2^{26}},
\]
which is a coefficient, not a variable. Both versions are integer polynomials. The difference of 17 variables comes from a chain specifying this fixed constant using smaller coefficients. The 850-variable version preserves the 909-variable baseline's maximum coefficient bit length of 3,247; its count is not asserted minimal among versions with unrestricted decimal coefficients. The new tables have 51,338 and 51,388 nonzero monomials, respectively, and their quadratic systems have 790 and 807 residuals.

The reduction from 909 to 850 comprises four changes: 36 variables from a new elliptic-curve chart and cancellation of ordinate products; nine variables from shorter Laurent-power circuits; one variable from scalar preparation; and 13 variables from exact, certified elimination. Removing the constant chain gives 833 variables, or 832 existential witnesses.

\begin{theorem}\label{thm:main}
Let $P_{909}$ be the direct-target polynomial generated by
\nolinkurl{research/overnight-20261106/unit_target.py}, variant~4, at the contribution-branch commit above. There are explicitly specified integer polynomials $P_{833}$ and $P_{850}$ of total degree four, in the indicated numbers of variables, such that for every $q\in\Q$,
\[
 \begin{aligned}
 &\exists\mathbf u\in\Q^{908}\;P_{909}(q,\mathbf u)=0\\
 \Longleftrightarrow\quad&\exists\mathbf v\in\Q^{832}\;P_{833}(q,\mathbf v)=0\\
 \Longleftrightarrow\quad&\exists\mathbf w\in\Q^{849}\;P_{850}(q,\mathbf w)=0.
 \end{aligned}
\]
The input coordinate is preserved pointwise. Consequently, if the source construction has rational zero projection $\Z$, so do both new polynomials.
\end{theorem}

The source's 909-variable polynomial has ordinary-table SHA-256
\begin{center}\small\texttt{326a23b5af8064915013411f24af39e4507c41431c6b783d5a6cae488695b93b}.\end{center}
The accompanying package contains exact residual systems, complete quartic tables, coordinate maps, elimination certificates and independent checking scripts. The following sections prove the nonmechanical chart replacement and explain the remaining transformations.

\section{An exact replacement of the point chart}
Put $K=\Q(\sqrt2)$, $\delta=3-\sqrt2$ and $D=\delta^2=11-6\sqrt2$. The old chart, with letters in $K$, is
\begin{equation}\label{eq:old}
 u=t^2,\qquad x=h^2,\qquad u(x-u/4)=D,\qquad y=h(u/2-x).
\end{equation}
The new chart is
\begin{equation}\label{eq:new}
 x=h^2,\qquad j^2=x^2-D,\qquad y=hj.
\end{equation}
An equation over $K$ means its two rational component equations. The last expression in \eqref{eq:new} need not be assigned its own coordinates.

\begin{lemma}\label{lem:squares}
If $U,V\in K$, $U\ne0$, and $V^2=U^3+4DU$, then $U$ is a square in $K$.
\end{lemma}
\begin{proof}
The ring $\mathcal O=\Z[\sqrt2]$ is a PID: rounding both coefficients of an element of $K$ to nearest integers leaves a remainder of absolute norm at most $1/2$, giving a Euclidean algorithm. Its units are $\pm(1+\sqrt2)^n$, and every totally positive unit is a square. For completeness, multiply a positive unit by a power of $1+\sqrt2$ to place its first real embedding in $[1,1+\sqrt2)$. Its other embedding is the positive or negative reciprocal. Writing the unit as $a+b\sqrt2$ then forces $a,b\in\{0,1\}$, and the norm condition and excluded upper endpoint leave only the unit~1. The two signs of the fundamental unit under the real embeddings give the assertion about totally positive units.

Let $m=v_{\mathfrak p}(U)$, with integral normalization. Away from $(\sqrt2)$ and $(\delta)$ the two terms $U^3$ and $4DU$ have valuations $3m$ and $m$. If $m$ is odd, their unequal minimum is odd, a contradiction. At $(\sqrt2)$ these valuations are $3m$ and $m+4$, with equality only at the even value $m=2$. Again $m$ cannot be odd. At $(\delta)$ they are $3m$ and $m+2$, with equality only at $m=1$. In that exceptional case write $U=\delta u$ with $u$ a local unit. Then
\[
 V^2=\delta^3u(u^2+4).
\]
The residue field is $\mathbf F_7$, in which $-4$ is a nonsquare. Thus the right side has valuation three, also impossible. Every finite valuation of $U$ is even.

Under each real embedding, $D>0$ and $U(U^2+4D)=V^2$, so $U>0$. By the PID property, $U=\gamma^2\eta$ for $\gamma\in K^*$ and a unit $\eta$. This unit is totally positive, hence a square.
\end{proof}

\begin{proposition}\label{prop:chart}
The charts \eqref{eq:old} and \eqref{eq:new} have exactly the same projected pairs $(x,y)\in K^2$. Every solution of the new chart has $h,j,x,y\ne0$, and $x$ is positive under both real embeddings.
\end{proposition}
\begin{proof}
Given \eqref{eq:old}, put $j=t^2/2-h^2$. Then
\[
 j^2=h^4-t^2h^2+t^4/4=h^4-D,
\]
and $y=hj$. Conversely, given \eqref{eq:new}, set $U=2(h^2+j)$ and $V=2hU$. The identity $j^2=h^4-D$ gives
\[
 U^2-4h^2U+4D=0,
\]
hence $V^2=U^3+4DU$. Also $U\ne0$, since $U=0$ would imply $D=0$. By Lemma~\ref{lem:squares}, choose $t\in K$ with $t^2=U$. It follows that
\[
 t^2(h^2-t^2/4)=D,\qquad h(t^2/2-h^2)=hj,
\]
as required.

If $h=0$, then $j^2=-D$, contradicting either real embedding. If $j=0$, then $h^4=D$, so $h^2=\pm\delta$. Taking norms yields $\N(h)^2=7$, impossible for $\N(h)\in\Q$. The remaining assertions follow immediately.
\end{proof}

The proof is specific to this field and this value of $D$. It does not use the generally false assertion that square abscissa characterizes doubles on every elliptic curve. It also avoids adding a Mordell--Weil rank or generator hypothesis to the chart replacement.

\section{Using the chart without ordinate variables}
All nine point modules use the new chart. It has six rational coordinates: the two components of each of $h,j,x$, subject to four quadratic component equations. The old module had ten coordinates, including the components of $t,h,u,x,y$. Thus the change saves four coordinates per point, provided each downstream ordinate is consumed as the expression $hj$.

For the difference of chart points $Q,R$, retain the existing inverse equations for $x_Q-x_R$ and define the slope by
\[
 (x_Q-x_R)m=h_Qj_Q+h_Rj_R.
\]
Its numerator and defining equation are quadratic in the available coordinates. Define the difference abscissa by the existing square of $m$, and its ordinate by
\[
 y_{Q-R}=m(x_Q-x_{Q-R})-h_Qj_Q.
\]
This again is a quadratic assignment. The retained inverse ensures that the slope has exactly its old value.

For a chart point,
\[
 z=-x/y=-h/j.
\]
Thus the local parameter is defined by $jz=-h$. When a logarithm value is combined with its parameter, replace $y\ell=-xG$ by $j\ell=-hG$, where $G$ is the unchanged truncated logarithm polynomial. At the scalar ratio replace
\[
 x_H\rho=-y_Hz_S\qquad\text{by}\qquad h_H\rho=-j_Hz_S.
\]
All cancellations are reversible by Proposition~\ref{prop:chart}. Difference-point denominators and all other guards remain unchanged. The fixed target has
\[
 h=-1+\tfrac74\sqrt2,\qquad j=\tfrac{39}{8}-\tfrac92\sqrt2;
\]
these give exactly its previous abscissa and ordinate.

Together with Proposition~\ref{prop:chart}, these identities give rational witness maps in both directions. To reconstruct the old chart, choose the square root supplied by Lemma~\ref{lem:squares}; then evaluate all old assignments. The choice of sign does not alter $(x,y)$ or any downstream value.

\section{Shorter weight circuits}
For each of nine detectors, write $e$ for its nonzero weight and $q$ for its order-one weight. The required outputs are $q_{12}=qe^{22}$ and $q_{40}=qe^{78}$. The old construction uses seven pure powers of $e$ and two final quotients. Instead construct
\[
 a_1=e^{-1},\quad a_3=a_1^2/e,\quad a_7=a_3^2/e,\quad
 a_{14}=a_7^2,\quad a_{29}=a_{14}^2/e,\quad a_{55}=a_{29}^2/a_3,
\]
and impose
\[
 a_{29}q_{12}=qa_7,\qquad a_{55}q_{40}=q_{12}e.
\]
These are six power coordinates and two output coordinates. They give the same outputs and save one coordinate per detector.

The nonzero condition is inherited, rather than presumed from a quotient notation. In the source, $w_0$ has a retained inverse and
\[
 w_0n=(w_0^2-2w_1^2)/2,\qquad (e\Delta)b=n^2.
\]
Since $w_0\ne0$ and 2 is a nonsquare in $\Q$, the first equation gives $n\ne0$. The second gives $e\Delta\ne0$. Therefore every Laurent denominator above is nonzero. Conversely all Laurent coordinates are uniquely defined for each source solution. The changes introduce no new zero branches and exclude no source solution.

\section{Scalar preparation with one fewer coordinate}
The scalar pole uses a nonzero input $b$ and a preparation witness $t$, with
\[
 U=\frac{b(1-bt^2)}{1+bt^2}.
\]
The old circuit stores $b^{-1},t_2=t^2,bt_2,U$. Replace these four coordinates by $t_2,T,V$, with
\[
 t_2=t^2,\qquad b(T-t_2)=1,\qquad TV=2,
\]
and use the affine expression $U=V-b$. The equation $b(T-t_2)=1$ forces $b\ne0$, and $TV=2$ forces $T\ne0$. Consequently $T=t^2+b^{-1}$, $1+bt^2=bT\ne0$, and
\[
 U=2/T-b=\frac{b(1-bt^2)}{1+bt^2}.
\]
Conversely every old solution extends uniquely by these values of $T,V$. Thus the modification saves one coordinate and preserves all downstream quantities. It also supplies the nonzero condition formerly imposed by the separate coordinate $b^{-1}$.

\section{Certified elimination and the final polynomial}
The ordinary compiler first performs its existing quadratic-preserving eliminations. The additional certificate removes 13 coordinates. Each step is either an invertible affine shear fixing the input, or elimination of a variable $v$ from an equation $cv-r=0$, where $c\in\Q^*$ and $r$ does not contain $v$.

After substituting $v=r/c$, an affected row may be adjusted by polynomial multiples of retained unaffected rows to cancel terms of degree greater than two. The independent checker verifies every such polynomial identity and requires all reducer rows to remain present. Thus the equations before and after adjustment have equivalent zero sets. In reverse, reinsert $v=r/c$ and invert the affine shears. No variable denominator is introduced.

Normalize each remaining residual by a nonzero rational constant to obtain integer quadratic polynomials $R_1,\ldots,R_m$, and set
\[
 P=\sum_{i=1}^mR_i^2.
\]
Over $\Q$, this vanishes exactly when all residuals vanish. The complete expansion is checked against this sum of squares. All surviving coordinates are active, and the input is fixed throughout. This proves Theorem~\ref{thm:main} once the preceding transformations are composed with the certified steps.

For the smaller-coefficient version, introduce 17 coordinates $c_{10},\ldots,c_{26}$ satisfying
\[
 c_{10}=3^{1024},\qquad c_{k+1}=c_k^2\quad(10\le k<26),
\]
and replace $S$ by $c_{26}$. These equations have a unique rational solution and keep every residual quadratic. The new polynomial therefore has the same projection and 17 additional variables.

\section{Verification scope and further results}
The supplied evidence separates the elementary chart proof, exact circuit identities, compiler replay, additional elimination replay, and full sum-of-squares expansion. Separate code checks the chart identities, all changed quotient formulas, the constant-chain specialization, coordinate support, degrees and coefficient integrality. These finite checks certify the transformations; they do not establish the analytic and arithmetic theorems inherited by the source construction.

In particular, the 909-variable baseline changes targets and comparison precisions relative to the main repository. Its arithmetic justification remains part of the inherited proof. The improvements in this note preserve its rational zero projection exactly, so they neither require a stronger prime-pattern input nor repair any gap that may remain in that source argument.

A separate variant starts directly from the main repository's 962-variable system and gives a quartic in 894 variables, keeping its original comparison precisions, target multiples and fixed coefficient $S=3^{2^{27}}$. The chart argument above applies; for target doublings, retain the target ordinates and replace the first occurrence of $x^2$ by $j^2+D$. An analogous Laurent chain, an exact normalization of the detector quotients, the scalar preparation and the same thirteen elimination types complete the reduction. The package includes its full residual system, polynomial, certificates and source. Thus this backup has precisely the main repository's rational zero projection, independently of the contribution branch's arithmetic refinements.

A separate constant-linear substitution search, beginning directly with the main repository's 962-variable system, gives a degree-six polynomial in 755 variables and a degree-eight polynomial in 601 variables. Their certificates were replayed independently. These are additional low-degree tradeoffs, not unrestricted variable records: the collaborator notes already contain much smaller constructions at very high degree.

\end{document}
